This paper develops a task-based framework for understanding how advances in AI reshape automation and production.
By modeling production as a sequence of steps that the firm aggregates into tasks, we show how firms optimally allocate work between humans and AI, whether in augmented or fully automated form.
A central feature of the model is the possibility of AI chains, in which multiple steps are executed by AI and only the final output is verified by a human.
We show that chaining can overturn comparative advantage logic in factor assignment by amplifying the importance of local work context.
That is, whether a unit of work is assigned to a human or to AI depends not only on the characteristics of the step in question but also on those of its neighbors in production.

We characterize the firm's optimal AI deployment decision across two horizons: a short run, in which the boundaries of jobs and what they pay are both fixed, so AI can only compress the time a workflow takes; and a long run, in which the firm redraws those boundaries jointly with its AI deployment, trading the lower wages of narrower jobs against the coordination costs of maintaining more of them.

We also show that improvements in AI quality can induce discrete reorganizations of work.
A key implication of chaining is that marginal increases in AI quality may generate little or no cost savings until a threshold is crossed, after which the optimal production configuration changes discontinuously, so that the marginal returns to improving AI are non-monotone.
Our framework thus helps explain why firms invest so heavily in AI capabilities even when short-run returns appear limited, and it is consistent with the J-curve pattern of technology adoption \citep{brynjolfsson2021productivity}, in which early adoption raises costs before the anticipated gains from reorganization are realized.

We complement our theoretical analysis with empirical evidence consistent with the core mechanisms implied by the model.
Using a task-level dataset of AI exposure and AI execution outcomes, we find that AI-executed steps tend to appear in contiguous AI chains, that occupations whose AI-exposed steps are more fragmented across the workflow exhibit lower realized AI execution conditional on exposure, and that when comparable steps appear across occupations, those adjacent to AI-executed neighbors are significantly more likely to be executed by AI themselves.
These patterns align with the central economic forces highlighted by the model: the importance of positional relationships among steps, the role of AI-able step dispersion in shaping execution outcomes, and the local complementarities created by AI chains.
